\section{Dataset and Data Structure}
Data is sourced from the \textbf{Clothing Dataset Full} (Kaggle). This dataset is selected because it closely simulates a real shop-owner workflow: each product record starts from essential fields (\texttt{image}, \texttt{label}, \texttt{color}), then the system automatically enriches metadata.

This design supports non-technical administrators who may not know prompt engineering. They only provide necessary product information, while the pipeline generates a clean, high-quality caption for retrieval and management.

\begin{table}[H]
\centering
\begin{tabular}{lll}
\toprule
\textbf{Column} & \textbf{Type} & \textbf{Description}\\
\midrule
image   & string & Image ID (filename without extension)\\
label   & string & Category label: T-Shirt, Dress, Pants, \ldots\\
color   & string & Primary color of the product\\
caption & string & Auto-enriched description generated by LLM (30--40 words)\\
\bottomrule
\end{tabular}
\caption{Fashion items data structure}
\end{table}

\section{Data Cleaning}

\begin{tcolorbox}[title=Method: Noise Label Filtering and Integrity Check]
Remove ambiguous labels (Not sure, Skip, Other) and verify image-file existence before ingestion.
\end{tcolorbox}

In the production source code (\texttt{pre\_processing/processing\_data.py}), noise label filtering is implemented via a constant set:

\begin{algorithm}[H]
\caption{Data Ingestion and Noise Filtering}\label{alg:data_filtering}
\begin{algorithmic}[1]
\Require CSV File $csv$, Image Directory $img\_dir$, Limit $L$
\Ensure Cleaned list of product tuples $R$
\Function{LoadAndFilterData}{$csv, img\_dir, L$}
    \State $R \gets []$
    \State $E \gets \{\text{``Not sure''}, \text{``Skip''}, \text{``Other''}\}$
    \For{each row in $csv$}
        \State $id \gets \text{row.image}$
        \State $label \gets \text{row.label}$
        \If{$id$ is empty \textbf{or} $label$ is empty \textbf{or} $label \in E$}
            \State \textbf{continue}
        \EndIf
        \State $path \gets img\_dir / \text{f``}\{id\}\text{.jpg''}$
        \If{\textbf{not} $path.\text{exists}()$}
            \State \textbf{continue}
        \EndIf
        \State Append $(id, label, path)$ to $R$
        \If{$\text{length of } R \ge L$}
            \State \textbf{break}
        \EndIf
    \EndFor
    \State \Return $R$
\EndFunction
\end{algorithmic}
\end{algorithm}

\textbf{Task breakdown:}
\begin{enumerate}
    \item \texttt{EXCLUDED\_LABELS} --- Set of labels with no clear fashion meaning.
    \item Verify image file existence before adding to the list.
    \item Return tuples \texttt{(image\_id, label, image\_path, source)} ready for PostgreSQL upsert.
\end{enumerate}

\section{Caption Generation via LLM (Gemini)}

This is a critical step to add \textbf{semantic tokens} to each product. Instead of relying on simple category labels (e.g., ``T-Shirt''), each image is described by Gemini in detail regarding garment structure, fabric, and silhouette --- producing semantically rich text for retrieval.

\begin{tcolorbox}[title=Method: Multimodal Caption Generation]
Gemini receives the actual product image as input and generates a 30--40 word description focusing on: (1) fabric/texture, (2) silhouette/fit, (3) construction details, (4) overall aesthetic. Colors and brand names are excluded to ensure generalizability.
\end{tcolorbox}

\begin{algorithm}[H]
\caption{Multimodal Caption and Color Enrichment via LLM}\label{alg:llm_enrichment}
\begin{algorithmic}[1]
\Require Image file $I$, Product Label $L$, LLM model $M$ (Gemini 2.5 Flash)
\Ensure Search Caption $C$, Primary Color $K$
\Function{GenerateSearchCaption}{$I$}
    \State $prompt \gets \text{``Write a 30-40 word fashion-centric morphological description...''}$
    \State $response \gets M.\text{generate\_content}(prompt, I)$
    \State $C \gets \text{Clean}(response.\text{text})$
    \State \Return $C$
\EndFunction
\State
\Function{DetectItemColor}{$I, L$}
    \State $prompt \gets \text{``Identify the primary color of this \{}L\text{\}...''}$
    \State $response \gets M.\text{generate\_content}(prompt, I)$
    \State $K \gets \text{Clean}(response.\text{text})$
    \State \Return $K$
\EndFunction
\end{algorithmic}
\end{algorithm}

\textbf{Task breakdown:}
\begin{itemize}
    \item \textbf{Multimodal input}: Gemini receives both a text prompt and a product image simultaneously, producing a semantic description.
    \item \textbf{Semantic token}: The resulting caption is a richer semantic token than a short label, e.g., \textit{``relaxed-fit woven cotton pullover with ribbed crew neckline and hemline, drop shoulders, minimal seaming''}.
    \item \textbf{Rate limiting}: The system processes in batches of 20 items with \texttt{asyncio.sleep()} to avoid rate limits.
    \item \textbf{Checkpoint}: PostgreSQL commits after each batch in case of disconnection.
\end{itemize}

\section{Color Detection via LLM}
Color enrichment is designed to produce compact but discriminative color tokens for retrieval. Instead of coarse labels (e.g., ``Green''), the model is prompted to return specific shades (e.g., ``Olive Green'', ``Charcoal Grey'') and pattern-aware output when relevant (e.g., ``White with Black stripes'').

\begin{tcolorbox}[title=Method: Color Enrichment for Retrieval]
For each product with a missing color, the pipeline loads the image, calls \texttt{detect\_item\_color()}, applies light output cleanup (\texttt{strip()} and markdown-symbol removal), and then upserts the result to \texttt{fashion\_item\_enrichment.color} with model metadata and timestamp fields.
\end{tcolorbox}

\textbf{Task breakdown:}
\begin{itemize}
    \item \textbf{Selective processing}: Only rows with \texttt{NULL}/empty color are selected (\texttt{fetch\_missing\_colors()}).
    \item \textbf{Prompt constraints}: Output must be color-only, specific, concise (max 3--4 words), and not full-sentence text.
    \item \textbf{Traceability}: Each item is logged as \texttt{success}, \texttt{failed}, or \texttt{skipped} in preprocessing logs.
    \item \textbf{Retrieval impact}: Enriched color tokens improve BM25 exact keyword matching on \texttt{label + color} and support downstream fuzzy soft relevance filtering.
\end{itemize}


% ===========================================================================
% CHAPTER 3: RAG v1.0 ARCHITECTURE
% ===========================================================================
